%!TEX root = ../main.tex

\subsection{Cone Decompositions}\label{app:cone-decomps}

\setlength{\tabcolsep}{4pt} % Default value: 6pt

\begin{table*}[t]
	\centering
	\caption{Approximating the C-ReLU problem with cone decompositions.
		We compare the solution to C-GReLU (FISTA) with cone-decomposition by solving the min-norm program (CD-SOCP),
		the approximate cone decomposition (CD-A), and directly solving C-ReLU using the AL method.
		We report the norm of each model to demonstrate the ``blow-up'' effect of the cone decomposition.
	}
	\label{table:cone-decomp-full}
	\vspace{0.1in}

	\begin{small}
		\begin{tabular}{lcccccccccccc}
			               & \multicolumn{3}{c}{\textbf{R-FISTA}} & \multicolumn{3}{c}{\textbf{CD-SOCP}} & \multicolumn{3}{c}{\textbf{CD-A}} & \multicolumn{3}{c}{\textbf{AL}}                                                                                                                                        \\ \cmidrule(lr){2-4}  \cmidrule(lr){5-7}  \cmidrule(lr){8-10} \cmidrule(lr){11-13}
			Dataset        & Acc.                                 & Time                                 & Norm                              & Acc.                            & Time   & Norm                          & Acc. & Time & Norm                           & Acc. & Time  & Norm                          \\ \midrule
			breast-cancer  & 68.4                                 & 0.84                                 & 1.06{\tiny \(\times 10^{2}\)}     & 68.4                            & 12.14  & 3.73{\tiny \(\times 10^{3}\)} & 68.4 & 3.18 & 4.76{\tiny \(\times 10^{3}\)}  & 66.7 & 6.17  & 5.84{\tiny \(\times 10^{1}\)} \\
			congressional  & 64.4                                 & 0.89                                 & 4.20{\tiny \(\times 10^{1}\)}     & 64.4                            & 80.69  & 1.90{\tiny \(\times 10^{3}\)} & 97.3 & 5.26 & 3.5{\tiny \(\times 10^{3}\)}   & 69.0 & 18.32 & 3.48{\tiny \(\times 10^{1}\)} \\
			sonar          & 87.8                                 & 0.45                                 & 2.25{\tiny \(\times 10^{1}\)}     & 90.2                            & 23.13  & 1.28{\tiny \(\times 10^{2}\)} & 64.4 & 1.09 & 1.83{\tiny \(\times 10^{2}\)}  & 87.8 & 4.34  & 2.33{\tiny \(\times 10^{1}\)} \\
			credit         & 82.6                                 & 0.44                                 & 7.32{\tiny \(\times 10^{1}\)}     & 83.3                            & 63.46  & 3.97{\tiny \(\times 10^{3}\)} & 84.1 & 5.1  & 4.76{\tiny \(\times 10^{3 }\)} & 84.1 & 6.46  & 5.11{\tiny \(\times 10^{1}\)} \\
			cylinder       & 75.5                                 & 1.18                                 & 9.64{\tiny \(\times 10^{1}\)}     & 77.5                            & 79.78  & 2.09{\tiny \(\times 10^{3}\)} & 75.5 & 2.68 & 2.61{\tiny \(\times 10^{3}\)}  & 75.5 & 15.18 & 1.11{\tiny \(\times 10^{2}\)} \\
			ecoli          & 71.6                                 & 0.07                                 & 1.73{\tiny \(\times 10^{1}\)}     & 71.6                            & 149.7  & 5.82{\tiny \(\times 10^{2}\)} & 70.1 & 0.36 & 1.13{\tiny \(\times 10^{2}\)}  & 70.1 & 3.38  & 1.68{\tiny \(\times 10^{1}\)} \\
			energy-y1      & 86.3                                 & 0.12                                 & 2.90{\tiny \(\times 10^{1}\)}     & 86.3                            & 134.55 & 2.34{\tiny \(\times 10^{3}\)} & 86.3 & 2.01 & 1.34{\tiny \(\times 10^{3}\)}  & 83.7 & 5.05  & 2.89{\tiny \(\times 10^{1}\)} \\
			glass          & 64.3                                 & 0.13                                 & 2.00{\tiny \(\times 10^{1}\)}     & 64.3                            & 68.76  & 4.05{\tiny \(\times 10^{2}\)} & 64.3 & 0.69 & 2.31{\tiny \(\times 10^{2}\)}  & 61.9 & 3.0   & 1.72{\tiny \(\times 10^{1}\)} \\
			cleveland      & 51.7                                 & 0.13                                 & 1.97{\tiny \(\times 10^{1}\)}     & 51.7                            & 109.4  & 2.52{\tiny \(\times 10^{2}\)} & 51.7 & 0.6  & 2.18{\tiny \(\times 10^{2 }\)} & 50.0 & 1.82  & 1.77{\tiny \(\times 10^{1}\)} \\
			hungarian      & 86.2                                 & 0.88                                 & 8.01{\tiny \(\times 10^{1}\)}     & 86.2                            & 17.15  & 2.75{\tiny \(\times 10^{3}\)} & 86.2 & 4.09 & 3.59{\tiny \(\times 10^{3}\)}  & 84.5 & 7.01  & 5.09{\tiny \(\times 10^{1}\)} \\
			heart-va       & 35.0                                 & 0.16                                 & 2.44{\tiny \(\times 10^{1}\)}     & 37.5                            & 67.11  & 4.79{\tiny \(\times 10^{2}\)} & 37.5 & 0.9  & 4.73{\tiny \(\times 10^{2 }\)} & 37.5 & 1.86  & 1.58{\tiny \(\times 10^{1}\)} \\
			hepatitis      & 80.6                                 & 1.06                                 & 2.96{\tiny \(\times 10^{1}\)}     & 80.6                            & 10.57  & 2.63{\tiny \(\times 10^{2}\)} & 80.6 & 1.97 & 3.4{\tiny \(\times 10^{2}\)}   & 74.2 & 8.56  & 5.53{\tiny \(\times 10^{1}\)} \\
			horse-colic    & 85.0                                 & 1.0                                  & 5.79{\tiny \(\times 10^{1}\)}     & 86.7                            & 29.42  & 9.45{\tiny \(\times 10^{2}\)} & 86.7 & 4.54 & 1.36{\tiny \(\times 10^{3}\)}  & 90.0 & 9.75  & 8.92{\tiny \(\times 10^{1}\)} \\
			ionosphere     & 90.0                                 & 1.15                                 & 4.74{\tiny \(\times 10^{1}\)}     & 90.0                            & 44.12  & 1.51{\tiny \(\times 10^{3}\)} & 90.0 & 5.76 & 2.05{\tiny \(\times 10^{3}\)}  & 90.0 & 15.91 & 6.42{\tiny \(\times 10^{1}\)} \\
			mammograph     & 78.1                                 & 0.24                                 & 4.13{\tiny \(\times 10^{1}\)}     & 78.1                            & 33.17  & 8.11{\tiny \(\times 10^{3}\)} & 78.1 & 5.18 & 3.02{\tiny \(\times 10^{3}\)}  & 79.2 & 5.14  & 3.70{\tiny \(\times 10^{1}\)} \\
			monks-2        & 60.6                                 & 1.95                                 & 1.02{\tiny \(\times 10^{2}\)}     & 60.6                            & 5.81   & 7.36{\tiny \(\times 10^{3}\)} & 57.6 & 6.61 & 9.43{\tiny \(\times 10^{3}\)}  & 45.5 & 18.31 & 8.13{\tiny \(\times 10^{1}\)} \\
			monks-3        & 87.5                                 & 1.6                                  & 5.59{\tiny \(\times 10^{1}\)}     & 87.5                            & 3.89   & 3.44{\tiny \(\times 10^{3}\)} & 87.5 & 6.49 & 4.49{\tiny \(\times 10^{3}\)}  & 95.8 & 31.84 & 8.60{\tiny \(\times 10^{1}\)} \\
			oocytes        & 78.6                                 & 0.98                                 & 1.14{\tiny \(\times 10^{2}\)}     & 79.1                            & 136.3  & 1.19{\tiny \(\times 10^{4}\)} & 78.0 & 5.67 & 1.3{\tiny \(\times 10^{4}\)}   & 74.2 & 81.68 & 8.54{\tiny \(\times 10^{1}\)} \\
			parkinsons     & 92.3                                 & 1.9                                  & 4.70{\tiny \(\times 10^{1}\)}     & 92.3                            & 16.03  & 1.06{\tiny \(\times 10^{3}\)} & 92.3 & 4.84 & 1.45{\tiny \(\times 10^{3}\)}  & 89.7 & 19.04 & 6.69{\tiny \(\times 10^{1}\)} \\
			pima           & 73.2                                 & 0.36                                 & 7.88{\tiny \(\times 10^{1}\)}     & 73.2                            & 37.68  & 8.09{\tiny \(\times 10^{3}\)} & 73.2 & 5.07 & 7.45{\tiny \(\times 10^{3}\)}  & 75.8 & 4.72  & 4.03{\tiny \(\times 10^{1}\)} \\
			planning       & 63.9                                 & 1.33                                 & 8.94{\tiny \(\times 10^{1}\)}     & 63.9                            & 10.54  & 2.16{\tiny \(\times 10^{3}\)} & 63.9 & 3.07 & 3.05{\tiny \(\times 10^{3}\)}  & 58.3 & 9.74  & 1.09{\tiny \(\times 10^{2}\)} \\
			seeds          & 95.2                                 & 0.25                                 & 1.91{\tiny \(\times 10^{1}\)}     & 95.2                            & 26.71  & 4.83{\tiny \(\times 10^{2}\)} & 95.2 & 1.99 & 4.17{\tiny \(\times 10^{2}\)}  & 95.2 & 5.16  & 1.82{\tiny \(\times 10^{1}\)} \\
			australian     & 64.5                                 & 0.69                                 & 1.58{\tiny \(\times 10^{2}\)}     & 63.8                            & 55.59  & 1.06{\tiny \(\times 10^{4}\)} & 64.5 & 5.46 & 1.21{\tiny \(\times 10^{4}\)}  & 65.2 & 4.8   & 2.74{\tiny \(\times 10^{1}\)} \\
			statlog-heart  & 81.5                                 & 0.81                                 & 6.76{\tiny \(\times 10^{1}\)}     & 81.5                            & 15.57  & 1.48{\tiny \(\times 10^{3}\)} & 81.5 & 2.5  & 2.02{\tiny \(\times 10^{3 }\)} & 85.2 & 7.36  & 6.11{\tiny \(\times 10^{1}\)} \\
			teaching       & 40.0                                 & 0.24                                 & 3.71{\tiny \(\times 10^{1}\)}     & 40.0                            & 12.18  & 1.10{\tiny \(\times 10^{3}\)} & 40.0 & 1.75 & 1.14{\tiny \(\times 10^{3}\)}  & 33.3 & 2.05  & 2.21{\tiny \(\times 10^{1}\)} \\
			tic-tac-toe    & 97.9                                 & 0.45                                 & 1.70{\tiny \(\times 10^{2}\)}     & 97.9                            & 60.36  & 6.86{\tiny \(\times 10^{3}\)} & 97.9 & 4.38 & 7.6{\tiny \(\times 10^{3}\)}   & 93.7 & 14.25 & 1.57{\tiny \(\times 10^{2}\)} \\
			vertebral-col. & 88.7                                 & 0.68                                 & 7.41{\tiny \(\times 10^{1}\)}     & 88.7                            & 9.64   & 7.76{\tiny \(\times 10^{3}\)} & 88.7 & 5.54 & 6.73{\tiny \(\times 10^{3}\)}  & 90.3 & 8.23  & 4.74{\tiny \(\times 10^{1}\)} \\
			wine           & 100                                  & 0.25                                 & 1.87{\tiny \(\times 10^{1}\)}     & 100                             & 32.05  & 1.97{\tiny \(\times 10^{2}\)} & 100  & 0.95 & 2.42{\tiny \(\times 10^{2}\)}  & 100  & 3.67  & 1.78{\tiny \(\times 10^{1}\)} \\ \bottomrule
		\end{tabular}
	\end{small}
\end{table*}

\setlength{\tabcolsep}{6pt}


Now we provide details and additional results for the cone-decomposition experiments given in \cref{table:cone-decomp}.

\textbf{Experimental Details}:
We selected 23 datasets from the UCI repository and fixed the regularization parameter at \( \lambda = 0.01 \).
Note that this parameter is not necessary optimal for each dataset;
the purpose of these experiments is to study the effects of using cone-decompositions to approximate the C-ReLU solution with a G-ReLU solution, rather than to obtain optimal test accuracies.
We randomly sampled 1000 activation patterns for the C-ReLU and C-GReLU models, removing duplicates and the zero pattern as necessary.
Note that we report the median results from five individual runs with re-sampled activation patterns to control for variance in the procedure.

For multi-class datasets, the convex formulations were extended as described in Appendix~\ref{app:multi-class-extension}.
We used the standard parameters for R-FISTA and the AL method as given in Appendix~\ref{app:default-parameters},
while the min-norm decomposition programs (CD-SOCP) was solved with MOSEK using the default parameters.
For CD-A, we set \( \lambda = 10^{-10} \) and \cref{eq:cd-approx} with R-FISTA
using the default parameters.
We terminate the optimization procedure when the min-norm subgradient has squared-norm
less than or equal to \( 10^{-10} \).
R-FISTA and the AL method were run on GPU compute nodes with one GeForce RTX 2080Ti graphics card, and four AMD 7502P CPUs, and 32 GB of RAM.
The cone decompositions were solved on identical nodes with four AMD 7502P CPUs with 32 GB of RAM.

\textbf{Additional Results}:
\cref{table:cone-decomp-full} provides the full set of results on all 23 datasets.
It also includes the final group norms of the models, calculated as
\[
	\sum_{D_i \in \tilde \calD} \norm{u_i^*}_2,
\]
for the C-GReLU model and
\[
	\sum_{D_i \in \tilde \calD} \norm{v_i^*}_2 + \norm{u_i^*}_2,
\]
for the C-ReLU models.
This allows us to quantify the ``blow-up'' in the model norm from decomposing \( u_i^* \) onto \( \calK_i - \calK_i \).
In practice, we find that CD-SOCP leads to very large increases in the model norm compared to the FISTA/AL solutions, while CD-A has a less severe effect.
However, the increased norms do not appear to affect the test accuracy of the final models.
Indeed, CD-SOCP and CD-A perform as well as the solution to the C-GReLU problem given by R-FISTA and are comparable to the AL method's solution.
The major downside of exact the cone-decomposition method is the huge increase in time necessary to solve for the decomposition.
This is largely because MOSEK is restricted to running on CPU.

